% !Mode:: "XeLaTeX:UTF-8"
%
% 内部书籍，谨慎参考！
% 撰写阶段错误难免~
%

%使用onesie去掉空白页
\documentclass[oneside, UTF8]{ctexbook}

%文献
\usepackage{natbib}

%章节标题左对齐和格式
\CTEXsetup[format={\LARGE\raggedright\bfseries}]{chapter} %修改每章的字体
\CTEXsetup[format={\Large\bfseries}]{section}
\CTEXsetup[format={\large\bfseries}]{subsection}
\CTEXsetup[format={\normalsize\bfseries}]{subsubsection}

%消除小字号警告
\usepackage{lmodern}
\usepackage{titling}

%页面设置
\usepackage[vcentering]{geometry}
\geometry{
	papersize={7in,9in},
	bottom=3pc,top=5pc,left=5pc,right=5pc,
	bmargin=4.5pc,
	footskip=18pt,
	headsep=25pt
}
\setlength\emergencystretch{1.5em}

%章节索引- 消除目录红框
\usepackage[pdfpagelabels=true,
pdffitwindow=false,
pdfview=FitH,
pdfstartview=FitH,
pagebackref=true,
breaklinks=true,
colorlinks=false,
bookmarks=true,
hidelinks,
bookmarksnumbered=true,
bookmarksopen=true,
bookmarksopenlevel=1,
bookmarksdepth=2,
plainpages=true]{hyperref}
\usepackage{bookmark}

%章节索引
%\usepackage[bookmarksnumbered]{hyperref}

%表格相关
\usepackage{multirow}
\usepackage{colortbl}
\usepackage{booktabs}

%数学公式
\usepackage{amsmath}
\usepackage{amsthm}
\usepackage{amssymb}
\usepackage{pgfplots}
\usepgfplotslibrary{colorbrewer}

%引用图片
\usepackage{graphicx}

%加水印
\usepackage{draftwatermark}
\SetWatermarkText{DRAFT}

%伪代码
\usepackage[chapter]{algorithm}
\usepackage{algorithmic}

%编号方式
\newcommand{\firstgls}[1]{\textbf{\,\gls{#1}}（\glsdesc{#1}）}
\newcommand{\firstacr}[1]{\textbf{\,\gls{#1}}（\glssymbol{#1}）}
\newcommand{\glsacr}[1]{\gls{#1}（\glssymbol{#1}）}
\newcommand{\firstall}[1]{\textbf{\,\gls{#1}}（\glsdesc{#1}, \glssymbol{#1}）}
\newcommand{\ENNAME}[1]{\text{#1}}
\newcommand{\NUMTEXT}[1]{\text{#1}}
\newcommand{\figref}[1]{图\,\ref{#1}\,}
\newcommand{\chapref}[1]{第\ref{#1}章}
\newcommand{\secref}[1]{第\,\ref{#1}\,节}
\newcommand{\algref}[1]{算法\,\ref{#1}\,}
\newcommand{\eqnref}[1]{式\,\eqref{#1}\,}
\newcommand{\tabref}[1]{表\,\ref{#1}\,}
\newcommand{\coderef}[1]{代码\,\ref{#1}\,}
\newcommand{\srclink}[1]{\underline{\color{blue}{#1}}}

%脚注改为①②-⑨样式
\usepackage{pifont}
\usepackage[perpage,symbol*]{footmisc}
\DefineFNsymbols{circled}{
     {\ding{192}}{\ding{193}}{\ding{194}}{\ding{195}}{\ding{196}}
     {\ding{197}}{\ding{198}}{\ding{199}}{\ding{200}}{\ding{201}}}
\setfnsymbol{circled}

% 支持表格分页 begin{longtable}, end{longtable}
\usepackage{longtable}

%流程图
\usepackage{tikz}
\usepackage{tikz-uml}
\usepackage{pgfopts}
\usepackage{calc}
\usepackage{xstring}
\usepackage{ifthen}
\usepackage{pgfkeys}
\usetikzlibrary{shapes.geometric, arrows, calc}
%流程图定义基本形状
%开始和结束
\tikzstyle{startstop} = [rectangle, rounded corners, minimum width = 1cm, minimum height=0.6cm,text centered, draw = black]
\tikzstyle{io} = [trapezium, trapezium left angle=70, trapezium right angle=110, minimum width=1cm, minimum height=0.6cm, text centered, draw=black]
%处理节点
\tikzstyle{process} = [rectangle, text centered, draw=black]
%条件判断
\tikzstyle{decision} = [diamond, aspect = 3, text centered, draw=black]
%runnable
\tikzstyle{runnable} = [circle, text centered, draw=black]
\tikzstyle{arrow} = [->,>=stealth] % 箭头形式
\tikzstyle{perceptron} = [circle,minimum size=2cm,draw=black]

%封面插入图片 ??暂时搞不清楚\\[\bigskipamount]双斜线用法
\preauthor{
\begin{center}
\includegraphics[width=6cm]{dl4j.png} \\[\bigskipamount]
}
\postauthor{\end{center}}

%标题间距
% From en book -B
\setlength{\parskip}{0.25 \baselineskip}

%新尺寸间距
% Sean said to make figures 26 picas wide
\newlength{\figwidth}
\setlength{\figwidth}{64pc}
% Spacing between notation sections
\newlength{\notationgap}
\setlength{\notationgap}{1pc}
% From en book -E

%支持段落首行不缩进\noindent
\usepackage{indentfirst}

%Java代码语法高亮
\usepackage{parcolumns} %分栏
\usepackage{listings}
\renewcommand{\lstlistingname}{代码}% Listing -> 代码
\usepackage{color}
\usepackage{xcolor}
\definecolor{dkgreen}{rgb}{0,0.6,0}
\definecolor{gray}{rgb}{0.5,0.5,0.5}
\definecolor{mauve}{rgb}{0.58,0,0.82}
\lstset{frame=tb,
     language=Java,
     aboveskip=3mm,
     belowskip=3mm,
     showstringspaces=false,
     columns=flexible,
     basicstyle = \ttfamily\small,
     numbers=left,
     numberstyle=\tiny\color{gray}, %设置行号的字体和颜色
     numbersep=4pt,  %设置行号与代码的距离，默认是5pt
     keywordstyle=\color{blue},
     commentstyle=\color{dkgreen},
     stringstyle=\color{mauve},
     breaklines=true,
     breakatwhitespace=true,
     tabsize=3
}

% http://tex.stackexchange.com/questions/198140/glossaries-and-custom-section-headings-broken  \glsentrytext!
\usepackage[nomain,acronym,xindy,toc,nopostdot]{glossaries}
\makeglossaries
\usepackage[xindy]{imakeidx}
\makeindex


% symbol and 
\include{math_symbol}
\include{terminology}

%文档开始
\begin{document}

%目录用罗马数字就行了
\frontmatter
%定义书籍标题
\title{\Huge\textbf{人工智能基础(Java)}}
\author{@YangJ9966, @imbbty, @simbaba}
\date{\today}

%显示书籍标题
\maketitle

%扉页
\centerline {\large {关于本书}}
\quad \par
2018年4月，教育部在研究制定《高等学校引领人工智能创新行动计划》，
并研究设立人工智能专业，进一步完善中国高校人工智能学科体系。

%书籍目录
\tableofcontents

%致谢
% !Mode:: "TeX:UTF-8"
\chapter*{致谢}
\addcontentsline{toc}{chapter}{致谢}

 我们才疏学浅而受此重任，深知自身水平难以将本书翻译得很准确。
因此我们完成草稿后，将书稿公开于Github，及早接受网友的批评和建议。
以下网友为本书的翻译草稿提供了很多及时的反馈和宝贵的修改意见：

 
所有校对的修改建议都保存在Github上，再次感谢以上同学和网友的付出。
经过这五个多月的修改，草稿慢慢变成了初稿。
尽管还有很多问题，但大部分内容是可读的，并且是准确的。
当然目前的翻译仍存在一些没有及时发现的问题，因此翻译也将持续更新，不断修改。
我们非常希望读者能到Github提建议，并且非常欢迎，无论多么小的修改建议，都是非常宝贵的。

此外，我们还要感谢XXX，帮助我们与出版社沟通交流，并给予了我们很多排版上的指导。

最后，感谢我们的杨教授，没有老师的支持，我们难以完成翻译。





\input{website}
\input{notation}

\mainmatter
% !Mode:: "TeX:UTF-8"
\part{综合运用}
\label{part:ai_exercise}

\newpage

%二维码
\begin{center}{\includegraphics[width=3cm]{barcode.png}}\end{center}\par

十年前，\emph{手写识别}还只是专业技术人员的玩物，而如今\emph{手写识别}已经变成深度学习的入门教程。
而\emph{模式识别}（Pattern Recognition）意味着需要专业人员提供很多特征。
模式识别从19世纪50年代开始出现，在20世纪80年代左右曾经风靡一时，也是信息科学和人工智能的重要组成部分。
它被应用于图像分析与处理、语音识别、数据挖掘等很多方面。
尽管模式识别也很高大上，也有了较长时间的应用，但效果一直差强人意。

这种由人提取特征交给机器，再交给机器去判断的方法识别率一直差强人意。
而\emph{机器学习}可以从海量数据学习\emph{特征}，使用已知经验数据（样本）自己提炼特征。
另外，再强调一下概念，深度学习只是机器学习算法的一种，而人工智能也不限于机器学习一种。

\begin{figure}[!htb] \centering
\begin{tikzpicture}
\node(learning) [draw, rectangle] at (2,0) {学习};
\node(knowledge) [draw, ellipse] at(2,1.5) {经验};
\draw[->](0,0) node[left]{新数据} -- node[above]{输入} (learning);
\draw[->](learning) -- node[above]{预测} (4,0) node[right]{人工智能};
\draw[->](knowledge) -- (learning);
\end{tikzpicture}
\end{figure}

大数据和深度学习，让我们能够摆脱人工特征提取的繁琐和束缚，让更多的人参与人工智能并为社会创造价值。
互联网+没有过时，再加上分布式和云计算，使得\emph{人工智能+}如鱼得水，再次迎来十年兴盛，相信接下来是数据为王的时代！

\chapter{大数据和分布式}
\label{chap:java_ex_cloud}

\section{云计算}

云计算是基于互联网的一种全新的模式。通过“云”我们可以轻松地存储数据，运行应用，
甚至进行强度很大的运算。云计算实质上实现了对计算机资源的完全掌控，
我们可以很轻易地通过云对资源进行分配，例如分配内存、处理器等计算资源。

机器学习和云计算、大数据的关系十分密切。大数据技术具有大量的数据与资源，能促进
机器学习的进一步发展。反过来，机器学习的发展也能促进数据挖掘，分析等方面
的进步。

\section{边缘计算}
边缘计算的定义是在靠近数据生成的设备端，就近提供服务。相比于云计算需要
将所有数据上传至云端的特性，边缘计算具有天然的优势。第一，由于边缘计算强调
“边缘”，数据的处理相比于云计算更为实时，高效；第二，边缘计算处理的数据多为
“小数据”，在数据存储和计算上成本都比较低；第三，边缘计算的出现降低了
对网络带宽的需求。随着物联网的不断发展，越来越多的联网设备将出现在
我们的生活中，这大大增加了网络传输的压力。边缘计算则通过数据本地处理
解决了这一问题。

边缘计算是对云计算的一种补充与完善。思科预计到2022年，移动设备和连接数量将达到123亿
\footnote{数据来源思科官网《思科移动网络VNI预测（2017-2022 年）》}。如果仅仅用
云计算来解决这么多互联设备的数据量往往会造成资源浪费，及时性较差，以及隐私泄露
等问题。边缘计算在一方面能减轻云端的负担，另一方面也能保护用户的隐私，具有重要
的实际意义。

\section{开源架构}
随着云计算和大数据的不断发展，人们对于数据存储和数据处理的效率
提出了越来越严格的要求，因此一批优秀的开源架构应运而生。

\subsection{Hadoop}
Hadoop是基于分布式计算的大数据开源架构，它的特点在于能高效，可靠地对
大量数据进行处理。Hadoop由HDFS和MapReduce组成，分别负责数据的分布式
储存和分布式计算。

HDFS提供了大规模存储的分布式解决方案，可以横跨成千上百台机器，操作却像对本地文件
进行操作一样简单，并能存储TB甚至PB级别的特大文件。这些优势来源于在HDFS中，
文件的储存是以较大的抽象块为单位，这样做的好处在于能将文件存储于
不同的硬盘中，使得文件的大小不受限于硬盘大小；同时较大的抽象块有利于对大型
文件进行寻址操作，最小化寻址开销。

除了抽象块，HDFS还包含了NameNode和DataNode。用户在用HDFS进行存储时，客户端首先
会将文件切块，再由DataNode储存在多个硬盘中。同时，NameNode负责记录每一个
文件的切块信息，以及每一块具体的存储机器。这样的系统就构成了NameNode为主服务器，多个
DataNode为从服务器的HDFS文件系统。

MapReduce是Hadoop的计算框架，能够以容错、可靠的方式并行处理海量的数据。
MapReduce主要包括了两个主要部分：Map阶段和Reduce阶段。



\subsection{Storm}
Storm是类似于Hadoop的实时数据处理框架，两者的区别在于Storm用于数据
的实时计算，Hadoop用于对离线数据进行处理。同时，Storm将数据存储在
内存中，这也有区别于Hadoop的HDFS存储方式。Storm也提供了类似于MapReduce
的简单编程模型，便于进行开发。

一个完整的Storm集群应该包括Nimbus，Zookeeper和Supervisor，工作架构
如图\ref{storm_framework}所示。其中，Nimbus负责接受客户端传来的拓扑代码
并分拆task，Zookeeper负责Nimbus和Supervisor之间的通信和协调，Supervisor
接收并处理数据。这样高效的工作架构使得Storm具有实时处理大批量数据的能力，
同时具有高扩展性，可以通过增加节点实现性能的线性提高。

~~~~
\begin{figure}[!ht]
\begin{center}
\begin{tikzpicture}
\tikzstyle{arrow} = [thick, -, >= stealth]
\tikzstyle{node} = [rectangle, rounded corners, minimum width = 2cm, minimum height=1cm,text centered, draw = black]
\node[node](Nimbus){Nimbus};
\node[node, right of = Nimbus, xshift = 2cm](zookeeper2){zookeeper};
\node[node, below of = zookeeper2, yshift = -0.5cm](zookeeper3){zookeeper};
\node[node, below of = zookeeper2, yshift = 2.5cm](zookeeper1){zookeeper};
\node[node, right of = zookeeper1, xshift = 2cm](supervisor1){supervisor};
\node[node, right of = zookeeper2, xshift = 2cm](supervisor2){supervisor};
\node[node, right of = zookeeper3, xshift = 2cm](supervisor3){supervisor};
\draw[arrow] (Nimbus)--(zookeeper1);
\draw[arrow] (Nimbus)--(zookeeper2);
\draw[arrow] (Nimbus)--(zookeeper3);
\draw[arrow] (zookeeper1)--(supervisor1);
\draw[arrow] (zookeeper1)--(supervisor2);
\draw[arrow] (zookeeper1)--(supervisor3);
\draw[arrow] (zookeeper2)--(supervisor1);
\draw[arrow] (zookeeper2)--(supervisor2);
\draw[arrow] (zookeeper2)--(supervisor3);
\draw[arrow] (zookeeper3)--(supervisor1);
\draw[arrow] (zookeeper3)--(supervisor2);
\draw[arrow] (zookeeper3)--(supervisor3);
\end{tikzpicture}
\caption{Storm架构示意图}
\label{storm_framework}
\end{center}
\end{figure}

\subsection{Spark}
Spark是一种通用的大数据计算框架，与Hadoop的MapReduce，Storm的
流式实时计算引擎十分相似。同时Spark也提供了大数据和机器学习方面的
框架，逐渐形成了大数据处理的一站式解决平台。

Spark的特点是支持大型，低延迟的数据处理方式。拓展了常用的MapReduce计算
模型。由于Spark支持在内存中进行计算，在必要时依赖硬盘进行复杂运算，因此
Spark比MapReduce更高效。

Spark是MapReduce的一种替代方案，兼容HDFS，Hive等存储方案，能够方便
地融入Hadoop的生态环境，以弥补MapReduce的不足。

\section{示例：人脸识别}



\chapter{车牌识别}
\label{chap:java_ex_recongnition}


手写识别机器学习领域中的一个经典问题。自从深度学习的兴起，MNIST已经是一个被”嚼烂”了的数据集。
该问题要解决的是把28x28像素的灰度手写数字图片，识别为0到9的数字。
而\emph{手写算式识别}在数字识别的基础上，需要识别运算符号，并输出计算结果。

\begin{figure}[!htb]
\centerline{\includegraphics[width=.1\figwidth]{images/handwriting_math.png} $\to$ 识别: $3+5 \times 6$} 
\label{fig:part4_handwriting_math}
\caption{手写算式识别}
\end{figure}

限于篇幅和笔者能力，本演示代码仅支持四则运算（$+ - \times \div$），其它运算符号留给同学们研究。
利用深度学习算法，为用户提供手写算式识别，可以极大的提高用户体验。
目前，有很多成熟的产品可以借鉴，相关的论文也非常多。

手写数字识别属于图像分类，可用的算法有：KNN、SVM、BP和CNN。
其中KNN和SVM比较善于数据挖掘，用做手写数字识别的话，比较难达到实用目的；
而BP使用全连接，难于构建深度网络，并存在梯度消失和过拟合问题。
本章例子，选用CNN深度学习训练手写算式识别。

\begin{figure}[!htb] \centering 
\begin{tikzpicture}
\node at (0,0) {\includegraphics[width=5cm,height=2cm]{images/handwriting_math.png}};
\draw (-2.5,1) rectangle (-1.6,-1);
\draw (-1.6,1) rectangle (-0.8,-1);
\draw (-0.8,1) rectangle (0.3,-1);
\draw (0.3,1) rectangle (1.3,-1);
\draw (1.3,1) rectangle (2.4,-1);
\end{tikzpicture}
\end{figure}

现在我们要识别一个算式，其实和MINIST数字识别没有太多差别，主要是增加了图像切割和四则运算符的识别任务。
为了降低难度，约束字符之间最好有一些间距，不要字符连笔。


\section{图像切割}
手写算式的输入的形式，有2种：拍照、手绘。照片要经过图像处理之后，才能进行切割；
而手绘可直接采集像素轨迹，就省掉了这一步。对于照片的预处理，需要先二值化和降噪把手写数字区域保留下来。
如果考虑图片的背景干扰，可先使用Canny算子检测边缘，这借助OpenCV的API很容易实现。

我们这种要求手写数字不黏连的情况，把Canny输出的图像往X轴投影，就能得出每个数字的横向宽度坐标。
虽然要求数字不黏连，但难保有横向重叠的情况，必须用类似\emph{范水法}进行切割。

\begin{figure}[!htb] \centering 
\begin{tikzpicture}
\node at (0,0) {\includegraphics[width=5cm]{images/handwriting_input.png}};
\node at (6,0) {\includegraphics[width=5cm]{images/handwriting_out.png}};
\end{tikzpicture}
\caption{图像处理：二值化和边缘识别(OpenCV)}
\end{figure}

\noindent
这种图像预处理，OpenCV\footnote{OpenCV是一个开源的跨平台视觉库，实现了图像处理和计算机视觉方面的很多通用算法，
支持Linux、Windows、Android和Mac OS操作系统上。
它轻量且高效，同时提供C++、Python、MATLAB等语言的接口。}非常擅长，不需要使用机器学习算法。


\noindent
\emph{Java代码：}
\begin{lstlisting}[language=Java]
System.loadLibrary(Core.NATIVE_LIBRARY_NAME);
Mat img = Imgcodecs.imread(src);
Imgproc.cvtColor(img, img, Imgproc.COLOR_BGR2GRAY);
Imgproc.Canny(img, img, threshold, threshold * 3, 3, true);
Imgcodecs.imwrite(dstImg, img);
\end{lstlisting}

\noindent
\emph{Python代码：}
\begin{lstlisting}[language=Python]
import cv2
from matplotlib import pyplot as plt

img = cv2.imread("handwriting_input.png", 0)
edges = cv2.Canny(img, 30, 60)
plt.subplot(1,2,1), plt.imshow(img, cmap='gray')
plt.title('Original Image'), plt.xticks([]), plt.yticks([])
plt.subplot(1,2,2), plt.imshow(edges, cmap='gray')
plt.title('Edge Image'), plt.xticks([]), plt.yticks([])
plt.show()
\end{lstlisting}

\noindent
实现Canny算子不在本书讨论范围，GitHub上的一份Java实现接近OpenCV的效果
\href{https://github.com/enzoftware/images-processing/blob/master/algorithms/Canny%20Filter/src/com/company/CannyFilter.java}
{\srclink{Github}上获取。}而CSDN上的一份实现简化了很多处理细节，有一定差距。有条件的同学，建议使用OpenCV的接口。
遇到的OpenCV问题，也容易在网络上检索。
\vspace{0.3cm}
\hrule
\begin{multicols}{2}
\begin{itemize}
\item[1.]转换为灰度图像
\item[2.]高斯模糊处理
\item[3.]根据梯度计算图像边缘幅值与角度
\item[4.]非最大信号压制处理（边缘细化）
\item[5.]双阈值边缘连接处理
\item[6.]二值化图像输出结果
\end{itemize}
\includegraphics[width=.12\figwidth]{images/Lenna.png}
\includegraphics[width=.12\figwidth]{images/Lenna_out.png}
% \end{tabular}
\end{multicols}
\hrule
\ \\
经过Canny算子处理之后，切割图像就非常容易了。把Canny输出的二值图像，在水平方向上投影，对数字进行竖向分割；
同理，在纵向再做一次，就能切割得紧凑一些。以下，是在原图上切割的效果 ：

\begin{figure}[!htb] \centering 
\begin{tikzpicture}
\node at (0,0) {\includegraphics[width=1cm]{images/3.png}};
\node at (2,0) {\includegraphics[width=1cm]{images/plus.png}};
\node at (4,0) {\includegraphics[width=1cm]{images/5.png}};
\node at (6,0) {\includegraphics[width=1cm]{images/times.png}};
\node at (8,0) {\includegraphics[width=1cm]{images/6.png}};
\end{tikzpicture}
\end{figure}

\noindent
Minist数据集是白底黑字的，使用它训练出来的模型，就很难识别我们的手写算式。所以，有必要在二值图像上切割。
\begin{figure}[!htb] \centering 
\begin{tikzpicture}
\node at (0,0) {\includegraphics[width=1cm]{images/3_c.png}};
\node at (2,0) {\includegraphics[width=1cm]{images/plus_c.png}};
\node at (4,0) {\includegraphics[width=1cm]{images/5_c.png}};
\node at (6,0) {\includegraphics[width=1cm]{images/times_c.png}};
\node at (8,0) {\includegraphics[width=1cm]{images/6_c.png}};
\end{tikzpicture}
\end{figure}

\noindent
但Canny生成的是空心字，并且有很多断线，使用FloodFill算法是填不满的。
笔者采用的是膨胀法，对于每一个\emph{白色点}向周围扩张6个像素，这样的结果还不错，差不多和Minist数据集很接近了。
\begin{figure}[!htb] \centering 
\begin{tikzpicture}
\node at (0,0) {\includegraphics[width=1cm]{images/3_b.png}};
\node at (2,0) {\includegraphics[width=1cm]{images/plus_b.png}};
\node at (4,0) {\includegraphics[width=1cm]{images/5_b.png}};
\node at (6,0) {\includegraphics[width=1cm]{images/times_b.png}};
\node at (8,0) {\includegraphics[width=1cm]{images/6_b.png}};
\end{tikzpicture}
\end{figure}

\noindent
尺寸也要标准化，做成和Minister一样的$28\times28$，这个对图片缩放就可以实现。


\section{算式识别}


\section{结果评估}
\chapter{写诗机器人}
\label{chap:java_ex_poem}

\section{word2vec介绍}
word2vec是google在2013年推出的一个NLP工具,它的特点是将所有的词向量化,这样词与词之间就可以定量的去度量他们之间的关系,挖掘词之间的联系。

\section{如何写诗}


 %JAVA
% !Mode:: "TeX:UTF-8"
\part{综合运用}
\label{part:ai_exercise}

\newpage

%二维码
\begin{center}{\includegraphics[width=3cm]{barcode.png}}\end{center}\par

十年前，\emph{手写识别}还只是专业技术人员的玩物，而如今\emph{手写识别}已经变成深度学习的入门教程。
而\emph{模式识别}（Pattern Recognition）意味着需要专业人员提供很多特征。
模式识别从19世纪50年代开始出现，在20世纪80年代左右曾经风靡一时，也是信息科学和人工智能的重要组成部分。
它被应用于图像分析与处理、语音识别、数据挖掘等很多方面。
尽管模式识别也很高大上，也有了较长时间的应用，但效果一直差强人意。

这种由人提取特征交给机器，再交给机器去判断的方法识别率一直差强人意。
而\emph{机器学习}可以从海量数据学习\emph{特征}，使用已知经验数据（样本）自己提炼特征。
另外，再强调一下概念，深度学习只是机器学习算法的一种，而人工智能也不限于机器学习一种。

\begin{figure}[!htb] \centering
\begin{tikzpicture}
\node(learning) [draw, rectangle] at (2,0) {学习};
\node(knowledge) [draw, ellipse] at(2,1.5) {经验};
\draw[->](0,0) node[left]{新数据} -- node[above]{输入} (learning);
\draw[->](learning) -- node[above]{预测} (4,0) node[right]{人工智能};
\draw[->](knowledge) -- (learning);
\end{tikzpicture}
\end{figure}

大数据和深度学习，让我们能够摆脱人工特征提取的繁琐和束缚，让更多的人参与人工智能并为社会创造价值。
互联网+没有过时，再加上分布式和云计算，使得\emph{人工智能+}如鱼得水，再次迎来十年兴盛，相信接下来是数据为王的时代！

\chapter{大数据和分布式}
\label{chap:java_ex_cloud}

\section{云计算}

云计算是基于互联网的一种全新的模式。通过“云”我们可以轻松地存储数据，运行应用，
甚至进行强度很大的运算。云计算实质上实现了对计算机资源的完全掌控，
我们可以很轻易地通过云对资源进行分配，例如分配内存、处理器等计算资源。

机器学习和云计算、大数据的关系十分密切。大数据技术具有大量的数据与资源，能促进
机器学习的进一步发展。反过来，机器学习的发展也能促进数据挖掘，分析等方面
的进步。

\section{边缘计算}
边缘计算的定义是在靠近数据生成的设备端，就近提供服务。相比于云计算需要
将所有数据上传至云端的特性，边缘计算具有天然的优势。第一，由于边缘计算强调
“边缘”，数据的处理相比于云计算更为实时，高效；第二，边缘计算处理的数据多为
“小数据”，在数据存储和计算上成本都比较低；第三，边缘计算的出现降低了
对网络带宽的需求。随着物联网的不断发展，越来越多的联网设备将出现在
我们的生活中，这大大增加了网络传输的压力。边缘计算则通过数据本地处理
解决了这一问题。

边缘计算是对云计算的一种补充与完善。思科预计到2022年，移动设备和连接数量将达到123亿
\footnote{数据来源思科官网《思科移动网络VNI预测（2017-2022 年）》}。如果仅仅用
云计算来解决这么多互联设备的数据量往往会造成资源浪费，及时性较差，以及隐私泄露
等问题。边缘计算在一方面能减轻云端的负担，另一方面也能保护用户的隐私，具有重要
的实际意义。

\section{开源架构}
随着云计算和大数据的不断发展，人们对于数据存储和数据处理的效率
提出了越来越严格的要求，因此一批优秀的开源架构应运而生。

\subsection{Hadoop}
Hadoop是基于分布式计算的大数据开源架构，它的特点在于能高效，可靠地对
大量数据进行处理。Hadoop由HDFS和MapReduce组成，分别负责数据的分布式
储存和分布式计算。

HDFS提供了大规模存储的分布式解决方案，可以横跨成千上百台机器，操作却像对本地文件
进行操作一样简单，并能存储TB甚至PB级别的特大文件。这些优势来源于在HDFS中，
文件的储存是以较大的抽象块为单位，这样做的好处在于能将文件存储于
不同的硬盘中，使得文件的大小不受限于硬盘大小；同时较大的抽象块有利于对大型
文件进行寻址操作，最小化寻址开销。

除了抽象块，HDFS还包含了NameNode和DataNode。用户在用HDFS进行存储时，客户端首先
会将文件切块，再由DataNode储存在多个硬盘中。同时，NameNode负责记录每一个
文件的切块信息，以及每一块具体的存储机器。这样的系统就构成了NameNode为主服务器，多个
DataNode为从服务器的HDFS文件系统。

MapReduce是Hadoop的计算框架，能够以容错、可靠的方式并行处理海量的数据。
MapReduce主要包括了两个主要部分：Map阶段和Reduce阶段。



\subsection{Storm}
Storm是类似于Hadoop的实时数据处理框架，两者的区别在于Storm用于数据
的实时计算，Hadoop用于对离线数据进行处理。同时，Storm将数据存储在
内存中，这也有区别于Hadoop的HDFS存储方式。Storm也提供了类似于MapReduce
的简单编程模型，便于进行开发。

一个完整的Storm集群应该包括Nimbus，Zookeeper和Supervisor，工作架构
如图\ref{storm_framework}所示。其中，Nimbus负责接受客户端传来的拓扑代码
并分拆task，Zookeeper负责Nimbus和Supervisor之间的通信和协调，Supervisor
接收并处理数据。这样高效的工作架构使得Storm具有实时处理大批量数据的能力，
同时具有高扩展性，可以通过增加节点实现性能的线性提高。

~~~~
\begin{figure}[!ht]
\begin{center}
\begin{tikzpicture}
\tikzstyle{arrow} = [thick, -, >= stealth]
\tikzstyle{node} = [rectangle, rounded corners, minimum width = 2cm, minimum height=1cm,text centered, draw = black]
\node[node](Nimbus){Nimbus};
\node[node, right of = Nimbus, xshift = 2cm](zookeeper2){zookeeper};
\node[node, below of = zookeeper2, yshift = -0.5cm](zookeeper3){zookeeper};
\node[node, below of = zookeeper2, yshift = 2.5cm](zookeeper1){zookeeper};
\node[node, right of = zookeeper1, xshift = 2cm](supervisor1){supervisor};
\node[node, right of = zookeeper2, xshift = 2cm](supervisor2){supervisor};
\node[node, right of = zookeeper3, xshift = 2cm](supervisor3){supervisor};
\draw[arrow] (Nimbus)--(zookeeper1);
\draw[arrow] (Nimbus)--(zookeeper2);
\draw[arrow] (Nimbus)--(zookeeper3);
\draw[arrow] (zookeeper1)--(supervisor1);
\draw[arrow] (zookeeper1)--(supervisor2);
\draw[arrow] (zookeeper1)--(supervisor3);
\draw[arrow] (zookeeper2)--(supervisor1);
\draw[arrow] (zookeeper2)--(supervisor2);
\draw[arrow] (zookeeper2)--(supervisor3);
\draw[arrow] (zookeeper3)--(supervisor1);
\draw[arrow] (zookeeper3)--(supervisor2);
\draw[arrow] (zookeeper3)--(supervisor3);
\end{tikzpicture}
\caption{Storm架构示意图}
\label{storm_framework}
\end{center}
\end{figure}

\subsection{Spark}
Spark是一种通用的大数据计算框架，与Hadoop的MapReduce，Storm的
流式实时计算引擎十分相似。同时Spark也提供了大数据和机器学习方面的
框架，逐渐形成了大数据处理的一站式解决平台。

Spark的特点是支持大型，低延迟的数据处理方式。拓展了常用的MapReduce计算
模型。由于Spark支持在内存中进行计算，在必要时依赖硬盘进行复杂运算，因此
Spark比MapReduce更高效。

Spark是MapReduce的一种替代方案，兼容HDFS，Hive等存储方案，能够方便
地融入Hadoop的生态环境，以弥补MapReduce的不足。

\section{示例：人脸识别}



\chapter{车牌识别}
\label{chap:java_ex_recongnition}


手写识别机器学习领域中的一个经典问题。自从深度学习的兴起，MNIST已经是一个被”嚼烂”了的数据集。
该问题要解决的是把28x28像素的灰度手写数字图片，识别为0到9的数字。
而\emph{手写算式识别}在数字识别的基础上，需要识别运算符号，并输出计算结果。

\begin{figure}[!htb]
\centerline{\includegraphics[width=.1\figwidth]{images/handwriting_math.png} $\to$ 识别: $3+5 \times 6$} 
\label{fig:part4_handwriting_math}
\caption{手写算式识别}
\end{figure}

限于篇幅和笔者能力，本演示代码仅支持四则运算（$+ - \times \div$），其它运算符号留给同学们研究。
利用深度学习算法，为用户提供手写算式识别，可以极大的提高用户体验。
目前，有很多成熟的产品可以借鉴，相关的论文也非常多。

手写数字识别属于图像分类，可用的算法有：KNN、SVM、BP和CNN。
其中KNN和SVM比较善于数据挖掘，用做手写数字识别的话，比较难达到实用目的；
而BP使用全连接，难于构建深度网络，并存在梯度消失和过拟合问题。
本章例子，选用CNN深度学习训练手写算式识别。

\begin{figure}[!htb] \centering 
\begin{tikzpicture}
\node at (0,0) {\includegraphics[width=5cm,height=2cm]{images/handwriting_math.png}};
\draw (-2.5,1) rectangle (-1.6,-1);
\draw (-1.6,1) rectangle (-0.8,-1);
\draw (-0.8,1) rectangle (0.3,-1);
\draw (0.3,1) rectangle (1.3,-1);
\draw (1.3,1) rectangle (2.4,-1);
\end{tikzpicture}
\end{figure}

现在我们要识别一个算式，其实和MINIST数字识别没有太多差别，主要是增加了图像切割和四则运算符的识别任务。
为了降低难度，约束字符之间最好有一些间距，不要字符连笔。


\section{图像切割}
手写算式的输入的形式，有2种：拍照、手绘。照片要经过图像处理之后，才能进行切割；
而手绘可直接采集像素轨迹，就省掉了这一步。对于照片的预处理，需要先二值化和降噪把手写数字区域保留下来。
如果考虑图片的背景干扰，可先使用Canny算子检测边缘，这借助OpenCV的API很容易实现。

我们这种要求手写数字不黏连的情况，把Canny输出的图像往X轴投影，就能得出每个数字的横向宽度坐标。
虽然要求数字不黏连，但难保有横向重叠的情况，必须用类似\emph{范水法}进行切割。

\begin{figure}[!htb] \centering 
\begin{tikzpicture}
\node at (0,0) {\includegraphics[width=5cm]{images/handwriting_input.png}};
\node at (6,0) {\includegraphics[width=5cm]{images/handwriting_out.png}};
\end{tikzpicture}
\caption{图像处理：二值化和边缘识别(OpenCV)}
\end{figure}

\noindent
这种图像预处理，OpenCV\footnote{OpenCV是一个开源的跨平台视觉库，实现了图像处理和计算机视觉方面的很多通用算法，
支持Linux、Windows、Android和Mac OS操作系统上。
它轻量且高效，同时提供C++、Python、MATLAB等语言的接口。}非常擅长，不需要使用机器学习算法。


\noindent
\emph{Java代码：}
\begin{lstlisting}[language=Java]
System.loadLibrary(Core.NATIVE_LIBRARY_NAME);
Mat img = Imgcodecs.imread(src);
Imgproc.cvtColor(img, img, Imgproc.COLOR_BGR2GRAY);
Imgproc.Canny(img, img, threshold, threshold * 3, 3, true);
Imgcodecs.imwrite(dstImg, img);
\end{lstlisting}

\noindent
\emph{Python代码：}
\begin{lstlisting}[language=Python]
import cv2
from matplotlib import pyplot as plt

img = cv2.imread("handwriting_input.png", 0)
edges = cv2.Canny(img, 30, 60)
plt.subplot(1,2,1), plt.imshow(img, cmap='gray')
plt.title('Original Image'), plt.xticks([]), plt.yticks([])
plt.subplot(1,2,2), plt.imshow(edges, cmap='gray')
plt.title('Edge Image'), plt.xticks([]), plt.yticks([])
plt.show()
\end{lstlisting}

\noindent
实现Canny算子不在本书讨论范围，GitHub上的一份Java实现接近OpenCV的效果
\href{https://github.com/enzoftware/images-processing/blob/master/algorithms/Canny%20Filter/src/com/company/CannyFilter.java}
{\srclink{Github}上获取。}而CSDN上的一份实现简化了很多处理细节，有一定差距。有条件的同学，建议使用OpenCV的接口。
遇到的OpenCV问题，也容易在网络上检索。
\vspace{0.3cm}
\hrule
\begin{multicols}{2}
\begin{itemize}
\item[1.]转换为灰度图像
\item[2.]高斯模糊处理
\item[3.]根据梯度计算图像边缘幅值与角度
\item[4.]非最大信号压制处理（边缘细化）
\item[5.]双阈值边缘连接处理
\item[6.]二值化图像输出结果
\end{itemize}
\includegraphics[width=.12\figwidth]{images/Lenna.png}
\includegraphics[width=.12\figwidth]{images/Lenna_out.png}
% \end{tabular}
\end{multicols}
\hrule
\ \\
经过Canny算子处理之后，切割图像就非常容易了。把Canny输出的二值图像，在水平方向上投影，对数字进行竖向分割；
同理，在纵向再做一次，就能切割得紧凑一些。以下，是在原图上切割的效果 ：

\begin{figure}[!htb] \centering 
\begin{tikzpicture}
\node at (0,0) {\includegraphics[width=1cm]{images/3.png}};
\node at (2,0) {\includegraphics[width=1cm]{images/plus.png}};
\node at (4,0) {\includegraphics[width=1cm]{images/5.png}};
\node at (6,0) {\includegraphics[width=1cm]{images/times.png}};
\node at (8,0) {\includegraphics[width=1cm]{images/6.png}};
\end{tikzpicture}
\end{figure}

\noindent
Minist数据集是白底黑字的，使用它训练出来的模型，就很难识别我们的手写算式。所以，有必要在二值图像上切割。
\begin{figure}[!htb] \centering 
\begin{tikzpicture}
\node at (0,0) {\includegraphics[width=1cm]{images/3_c.png}};
\node at (2,0) {\includegraphics[width=1cm]{images/plus_c.png}};
\node at (4,0) {\includegraphics[width=1cm]{images/5_c.png}};
\node at (6,0) {\includegraphics[width=1cm]{images/times_c.png}};
\node at (8,0) {\includegraphics[width=1cm]{images/6_c.png}};
\end{tikzpicture}
\end{figure}

\noindent
但Canny生成的是空心字，并且有很多断线，使用FloodFill算法是填不满的。
笔者采用的是膨胀法，对于每一个\emph{白色点}向周围扩张6个像素，这样的结果还不错，差不多和Minist数据集很接近了。
\begin{figure}[!htb] \centering 
\begin{tikzpicture}
\node at (0,0) {\includegraphics[width=1cm]{images/3_b.png}};
\node at (2,0) {\includegraphics[width=1cm]{images/plus_b.png}};
\node at (4,0) {\includegraphics[width=1cm]{images/5_b.png}};
\node at (6,0) {\includegraphics[width=1cm]{images/times_b.png}};
\node at (8,0) {\includegraphics[width=1cm]{images/6_b.png}};
\end{tikzpicture}
\end{figure}

\noindent
尺寸也要标准化，做成和Minister一样的$28\times28$，这个对图片缩放就可以实现。


\section{算式识别}


\section{结果评估}
\chapter{写诗机器人}
\label{chap:java_ex_poem}

\section{word2vec介绍}
word2vec是google在2013年推出的一个NLP工具,它的特点是将所有的词向量化,这样词与词之间就可以定量的去度量他们之间的关系,挖掘词之间的联系。

\section{如何写诗}


 %MATH
% !Mode:: "TeX:UTF-8"
\part{人工智能应用}
\label{part:ai_learning}

\newpage

%二维码
\begin{center}{\includegraphics[width=3cm]{barcode.png}}\end{center}\par

在\gls{AI}的早期，那些对人类智力来说非常困难、但对计算机来说相对简单的问题得到迅速解决，比如，那些可以通过一系列形式化的数学规则来描述的问题。
\gls{AI}的真正挑战在于解决那些对人来说很容易执行、但很难形式化描述的任务，如识别人们所说的话或图像中的脸。对于这些问题，我们人类往往可以凭借直觉轻易地解决。

线性回归是机器学习中最基本的一个算法。回归分析中，又依据描述自变量与因变量之间因果关系的函数表达式是线性的还是非线性的，分为线性回归分析和非线性回归分析。

\chapter{ND4J}
\label{chap:java_nd4j}

\section{ND4J介绍}

Nd4j是一个基于JVM的,接口和numpy接近的张量运算库。目前流行的数据分析框架，大多是基于Python或MatLab开发的，很难移植到JVM运行环境。
虽然有Colt这样的Java类库，但它对商业不友好。
\vspace{0.5cm}\noindent

\emph{ND4J主要特点：}（摘自官网）
\begin{itemize}
\item[1.] 多用途多维数组对象
\item[2.] 支持多平台扩展，包括GPU
\item[3.] 提供线性代数和信号处理功能
\end{itemize}

\noindent
使用ND4J创建多维向量和矩阵非常容易，用起来和numpy非常类似，甚至还提供了ND4S这种更加接近numpy的类库。

\noindent \ \\
\emph{创建2X2的NDArray：}
\begin{lstlisting}[language=Java]
INDArray arr1 = Nd4j.create(new float[]{1,2,3,4},new int[]{2,2});
System.out.println(arr1);
\end{lstlisting}
输出：
[[1.0 ,3.0]
[2.0 ,4.0]
]

\section{使用ND4J}
在学习使用ND4J之前，有必要了解数组在内存中的组织方式。数组通常是一段连续的内存，但Java并不能保证整个维度是连续的，这也很容易证明：

\begin{lstlisting}[language=Java,caption={Java二维数组},label=code:part3_array_mem]
int[][] numbers = new int[2][];
int[] N1 = {1,2,3};
int[] N2 = {2,3,4};

numbers[0] = N1;
numbers[1] = N2;
\end{lstlisting}

\coderef{code:part3_array_mem}，展示了N1和N2是2个不同的int[]对象，但无法确定数据在内存中的组织是否连续。Java这种分散的数据组织范式，非常不利于批量运算和GPU加速。
为了更有效率的使用数据，ND4J没采用Java的多维数组形式来表示张量（标量、向量、矩阵），而是在JVM之外申请的内存。
不仅有更好的性能，还可以结合更多的Native加速线性代数库，譬如BLAS。
\footnote{BLAS：Basic Linear Algebra Subprograms，目前最快的是Intel的MKL}

\begin{figure}[!htb]
\centerline{\includegraphics[width=.35\figwidth]{images/c_vs_f_order.png}}
\caption{C序和F序}
\label{fig:part3_c_vs_f_order}
\end{figure}

\noindent
\figref{fig:part3_c_vs_f_order}展示了这种不同。创建NDArray的时候，可以指定使用哪种方式。

\subsection{创建和更新}
与前面IntVector不同的是，NDArray的数据只有double或float类型，并且所有的NDArray共享相同的dtype。默认情况下，NDArray的数据都是float类型的。
创建NDArray有很多方式，ND4J还为全0和全1数组提供了便利接口：
\vspace{0.3cm}

\begin{lstlisting}[language=Java]
Nd4j.zeros(int...)
Nd4j.ones(int...)
\end{lstlisting}
\noindent
其中，int...是维度信息。

创建一般的NDArray则可以用章节一开始时提到的create()方法，传入两个数组作为数据组和维度即可。

另外，对于元素服从正态分布的NDArray，或是元素属于(0,1)区间的NDArray，ND4J都提供了相应的接口：
\begin{lstlisting}[language=Java]
Nd4j.rand(int,int);
Nd4j.randn(int,int);
\end{lstlisting}

在DL4J中，更新操作也很便利，利用putScalar和putRow方法，传入需要改变元素的位置与值，即可完成操作，
例如下面的代码，putScalar方法将矩阵下标为(1,3)的元素修改成100，putRow方法则将第一行
的所有元素重新设置成传入的INDArray中对应的元素，这样就完成了元素的更新。

\begin{lstlisting}[language=Java]
arr1.putScalar(1, 3, 100);
arr1.putRow(1,INDArray);
\end{lstlisting}

\subsection{切片操作}

对于NDArray，我们往往需要其中的一部分来进行操作，这就需要用切片操作来完成。
切片操作的代码如下图所示，我们传入一个或多个行数信息，就能获得对应行数的NDArray。
但是需要注意的是这两种方法只适用于二维矩阵。

\begin{lstlisting}[language=Java]
arr1.getRow(int);
arr1.getRows(int...);
\end{lstlisting}


但在日常使用中，我们对NDArray的切片要求往往更加复杂，仅仅使用上述的两个方法会比较困难。
例如我们对一个NDArray进行取前三行，前三列的操作时，仅仅用getRow和getRows方法就显得比较麻烦。
这时我们就需要用get方法进行操作：

\begin{lstlisting}[language=Java]
arr1.create(new float[]{1,2,3,4,5,6,7,8,9},new int[]{3,3});
INDArray arr2 = arr1.get(NDArrayIndex.interval(0,2), NDArrayIndex.interval(0,2));
System.out.println(arr2);
\end{lstlisting}

\noindent
输出结果为[[1.0,2.0] [4.0,5.0]]，这就完成了对NDArray的切片操作。

\subsection{基本运算}

NDArray的基本运算包括张量操作和元素操作。其中张量操作包括加减乘除四则运算，
操作对NDArray中的所有元素进行，在调用时传入一个数即可，例如：

\begin{lstlisting}[language=Java]
doube myDouble = 1.0;
INDArray arr1 = Nd4j.create(new float[]{1,2},new int[]{1, 2}); 
NDArray arr2 = arr1.add(myDouble)；
arr1.sub(myDouble)；
arr1.mul(myDouble)；
arr1.add(myDouble)；
System.out.println(arr2);
\end{lstlisting}

\noindent
输出结果为[[2.0,3.0]]。


元素操作与张量操作类似，都是对NDArray中的元素进行运算操作，不同的是张量操作传入的是一个数，
所有元素执行同一个操作；而元素操作传入一个NDArray，每个元素与另一个NDArray中对应的元素进行
操作。因此，元素操作要求进行操作的两个矩阵形状一致。

\begin{lstlisting}[language=Java]
INDArray arr1 = Nd4j.create(new float[]{1,2},new int[]{1, 2});
INDArray arr2 = Nd4j.create(new float[]{2,3},new int[]{1, 2});
System.out.println(arr1.add(arr2));
\end{lstlisting}

\noindent
输出结果为[[3.0,5.0]]。


\section{线性回归示例}









\chapter{DL4J}
\label{chap:java_dl4j}

\section{DL4J介绍}

Deeplearning4j(简称DL4J)是为Java和Scala编写的首个商业级开源分布式深度学习库。
DL4J还可与Hadoop和Spark集成,为商业化提供了很好的支持，而不仅仅是学术研究工具。
而在使用Deep Learning之前，我们还要先确定问题的类型，是有监督学习还是无监督问题？
对于有监督学习，你需要先提供有标签的数据进行训练；对于无监督学习，需要能自主检测数据的相似性和异常状况。

还要考虑处理的特征数量有多少？特征数量越多，需要的内存也越大。就图像而言，第一层的特征数量等于图像所包含的像素数。
所以MNIST数据集中的28 x 28像素的图像有784个特征。医疗诊断中的图像则可能有14兆像素。

\ \\
\begin{figure}[!htb]
\centerline{\includegraphics{images/dl4j_logo.png}}
\label{fig:part3_dl4j_logo}
\end{figure}

主流的深度学习框架有TensorFlow、Caffe、Keras、Theano等，它们大多在单节点服务器通过GPU加速完成模型训练。
随着大数据时代的来临，使用分布式计算极大提高了计算性能。因此将分布式计算与深度学习结合成为必然趋势。
DL4J就是为此而生，利用Spark在多服务器多GPU上进行分布式的深度学习模型训练，让模型跑得更快。
运用DL4J高效的训练一个完整神经网络模型包括：
\begin{itemize}
\item[1.] 预加载数据，对数据进行预处理。DL4J提供了\emph{DataVec}简化数据采集过程。
\item[2.] 超参数配置，各种优化算法和学习率、激活函数等。
\item[3.] 构建神经网络，从单层到多层神经网络的构建。
\item[4.] 模型训练，借助ND4J和分布式计算，高效地训练模型。
\item[5.] 模型评估和保存。这正是DL4J的核心算法部分。
\end{itemize}

所以，在开始训练之前，要先花点时间学习如何使用DL4J的构建神经网络以及加载数据。
由于加载数据的过程（reshape、正则化、归一化等）很容易引入意想不到的错误，
DL4J提供了\emph{DataVec}类对数据加工预处理，不仅可以减少错误，还支持多种数据源，譬如图片、CSV、ARFF和纯文本等。

\section{构造网络}
DL4J的神经网络都是用NeuralNetConfiguration构建的，用它可以设定网络的层（layer）以及超参数。
典型的神经网络包含：输入和输出层，以及隐藏层。对于全连接层，在DL4J中也称为\emph{DenseLayer}。\\

\begin{figure}[!htb]
\centerline{\includegraphics[width=.35\figwidth]{images/onelayer_labeled.png}}
\label{fig:part3_dl4j_onelayer}
\end{figure}


1. 设置超参数
\begin{lstlisting}[language=Java]
MultiLayerConfiguration conf = new NeuralNetConfiguration.Builder()
            .seed(rngSeed)
            .optimizationAlgo(OptimizationAlgorithm.STOCHASTIC_GRADIENT_DESCENT)
            .iterations(1)
            .learningRate(0.006)
            .updater(Updater.NESTEROVS).momentum(0.9)
            .regularization(true).l2(1e-4)
            .list()
\end{lstlisting}
\emph{超参数}(Hyperparameter)决定了神经网络学习方式的变量，包括模型的权重、激活函数、学习率和优化算法等。
\vspace{0.3cm}

2. 构建神经网络层
\begin{lstlisting}[language=Java]
.layer(0, new DenseLayer.Builder()
            .nIn(numRows * numColumns) // Number of input datapoints.
            .nOut(1000) // Number of output datapoints.
            .activation("relu") // Activation function.
            .weightInit(WeightInit.XAVIER) // Weight initialization.
            .build())
.layer(1, new OutputLayer.Builder(LossFunction.NEGATIVELOGLIKELIHOOD)
            .nIn(1000)
            .nOut(outputNum)
            .activation("softmax")
            .weightInit(WeightInit.XAVIER)
            .build())
.pretrain(false).backprop(true)
.build();
\end{lstlisting}

而\emph{超参数}(Hyperparameter)决定了神经网络学习方式的变量，包括模型的权重、激活函数、学习率和优化算法等。



\section{准备数据}
再好的训练模型，如果没有数据支撑，得出的结果也不是你想要的。本节将详细介绍常规数据使用方法，以及规范化手段。

\subsection{MNIST数据集}
MNIST（Mixed National Institute of Standards and Technology database）
是一个计算机视觉数据集，包含60000张手写数字的灰度图片训练集(training set) ，以及10000图片测试集(test set) 。
由来自250个不同人手写的数字构成，其中 50\% 是高中学生, 50\% 来自人口普查局 (the Census Bureau) 的工作人员。
测试集也是同样比例的手写数字数据。其中每一张图片都包含 $28 \times 28$ 个像素点，并配有对应的标签，也即是图片对应的数字。
\vspace{0.6cm}

\noindent
\figref{fig:part3_mnist_dataset}，展示了测试集的组织形式：

\begin{figure}[!htb]
\centerline{\includegraphics[width=.35\figwidth]{images/MNIST-dataset.png}}
\label{fig:part3_mnist_dataset}
\caption{MNIST数据}
\end{figure}

\noindent
MNIST官网提供了以下4个下载链接：
\begin{itemize}
\item[1.] train-images-idx3-ubyte.gz:  training set images (9912422 bytes) 
\item[2.] train-labels-idx1-ubyte.gz:  training set labels (28881 bytes) 
\item[3.] t10k-images-idx3-ubyte.gz:   test set images (1648877 bytes) 
\item[4.] t10k-labels-idx1-ubyte.gz:   test set labels (4542 bytes)
\end{itemize}

\noindent
看到这里，大家最更关心的还是下载的MNIST数据集，是如何定义的。这可以在官网查到，格式如下：
\vspace{0.3cm}

\noindent \textbf{训练集标签文件(train-labels-idx1-ubyte):}
\begin{lstlisting}
[偏移]    [数据类型]             [值]              [描述] 
0000       32 bit integer       0x00000801(2049)     magic number (MSB first) 
0004       32 bit integer       60000                number of items 
0008       unsigned byte        ??                   label 
0009       unsigned byte        ??                   label 
........ 
xxxx       unsigned byte        ??                   label
标签的值：0 ~ 9

16进制数据： 训练集-标签集合
0000 0801 0000 ea60 0500 0401 0902 0103
0104 0305 0306 0107 0208 0609 0400 0901
0102 0403 0207 0308 0609 0005 0600 0706
.....
\end{lstlisting}

\noindent \textbf{训练集数据文件(train-images-idx3-ubyte):}
\begin{lstlisting}
[偏移]     [数据类型]             [值]             [描述] 
0000       32 bit integer       0x00000803(2051)     magic number 
0004       32 bit integer       60000                number of images 
0008       32 bit integer       28                   number of rows 
0012       32 bit integer       28                   number of columns 
0016       unsigned byte        ??                   pixel 
0017       unsigned byte        ??                   pixel 
........ 
xxxx       unsigned byte        ??                   pixel

16进制数据： 训练集-图片集合
0000 0803 0000 ea60 0000 001c 0000 001c
0000 0000 0000 0000 0000 0000 0000 0000
.....
0000 0000 0000 0000 0000 0000 0000 0000
0000 0000 0000 0000 0312 1212 7e88 af1a
a6ff f77f 0000 0000 0000 0000 0000 0000
1e24 5e9a aafd fdfd fdfd e1ac fdf2 c340
.....
\end{lstlisting}

DL4J对MNIST数据集支持得很好，提供有deeplearning4j-datasets-xxxx.jar，其中有很多便利的API接口。



\subsection{IRIS数据集}
Iris也称鸢尾花卉数据集，是常用的分类实验数据集，由Fisher, 1936收集整理。包含150个数据集，分为3类，每类50个数据，每个数据包含4个属性。
可通过花萼长度，花萼宽度，花瓣长度，花瓣宽度4个属性预测鸢尾花卉属于（Setosa，Versicolour，Virginica）三个种类中的哪一类。

\begin{figure}[!htb]
\centerline{\includegraphics[width=.4\figwidth]{images/iris.png}}
\label{fig:part3_iris_flower}
\caption{莺尾花}
\end{figure}

该数据集由3种不同类型的鸢尾花的50个样本数据构成。其中的一个种类与另外两个种类是线性可分离的，后两个种类是非线性可分离的。
Iris以鸢尾花的特征作为数据来源，常用在分类操作中。该数据集包含了5个属性：

\begin{itemize}
\item[1.] Sepal.Length（花萼长度），单位是cm;
\item[2.] Sepal.Width（花萼宽度），单位是cm;
\item[3.] Petal.Length（花瓣长度），单位是cm;
\item[4.] Petal.Width（花瓣宽度），单位是cm;
\item[5.] 种类：Iris Setosa（山鸢尾）、Iris Versicolour（杂色鸢尾），以及Iris Virginica（维吉尼亚鸢尾）。
\end{itemize}


\subsection{波斯顿房价数据集}
波士顿房价数据集（Boston House Price Dataset）
，该数据集是美国人口普查局收集的有关波士顿马萨诸塞州住房的信息，发表于1978年美国某经济学杂志。
该数据集包含波士顿房屋的价格及其各项数据，每个数据项有14个数据，常用于回归分析。
\begin{itemize}
\item CRIM    - 城镇人均犯罪率
\item ZN      - 占地面积超过25,000平方英尺的住宅用地比例。
\item INDUS   - 城镇非零售商用土地的比例。
\item CHAS    - Charles River虚拟变量（如果是河道，则为1;否则为0）。
\item NOX     - 一氧化氮浓度（每千万）
\item RM      - 住宅平均房间数。
\item 年龄     - 1940年之前建成的自用房屋比例。
\item DIS     - 到波士顿五个中心区域的加权距离。
\item RAD     - 径向高速公路的可达性指数
\item 税       - 每10,000美元的全额物业税率
\item PTRATIO - 城镇的学生与教师比例
\item B       - $1000（Bk - 0.63）^ 2$，其中Bk是城镇黑人的比例
\item LSTAT   - 人口中地位低下者的比例。
\item MEDV    - 自有住房的中位数价值1000美元
\end{itemize}

波士顿房价数据集，共有506个观察，有13个输入变量和1个输出变量。在回归分析中，如果有两个或两个以上的自变量，就称为多元回归。
使用多个自变量的共同来预测，比只用一个自变量进行预测更符合实际，因此多元线性回归比一元线性回归的实用意义更大。
实际上，经常抽取波士顿房价数据的1-2个特征，做线性回归分析的演示。

\section{配置DL4J}
使用InteliJ IDEA开发环境，建议创建Maven工程，根据官网的说明配置pom.xml联网下载依赖就可以。
如果你比较熟悉Gradle，官网也有相应的说明。

\begin{lstlisting}[language=XML, caption={pom.xml}]
<dependencies>
  <dependency>
      <groupId>org.deeplearning4j</groupId>
      <artifactId>deeplearning4j-core</artifactId>
      <version>1.0.0-beta2</version>
  </dependency>
  <dependency>
      <groupId>org.nd4j</groupId>
      <artifactId>nd4j-native-plaform</artifactId>
      <version>1.0.0-beta2</version>
  </dependency>
</dependencies>
\end{lstlisting}


\section{线性回归}
线性回归可分为\emph{一元线性回归}和\emph{多元线性回归}。
只包括一个自变量和一个因变量的线性回归，结果近似一条直线。
\vspace{0.3cm}

\begin{tikzpicture}[scale=1.2]
    % Draw axes
    \draw [<->,thick] (0,2.5) node (yaxis) [above] {$y$}
        |- (4,0) node (xaxis) [right] {$x$};
    
    \draw (0,0.3) coordinate (a_1) -- (3,2.4) coordinate (a_2);

    \fill[red] (0.2,0.3) circle (2pt);
    \fill[red] (0.4,0.6) circle (2pt);
    \fill[red] (0.6,1.0) circle (2pt);

    \fill[red] (1.0,1.2) circle (2pt);
    \fill[red] (1.4,1.2) circle (2pt);
    \fill[red] (1.6,1.7) circle (2pt);

    \fill[red] (2.0,1.5) circle (2pt);
    \fill[red] (2.3,1.8) circle (2pt);

    \node[right=2.0cm, above=1.6cm] at (xaxis) {y=ax+b，其中x和y是已知的。};
    \node[right=2.0cm, above=0.6cm] at (xaxis) {求a=?,b=?};

    \draw[red,dashed] plot[smooth] coordinates {
        (0.2,0.3) (0.4,0.6) (0.6,1.0) (1.0,1.2)
        (1.4,1.2) (1.6,1.7) (2.0,1.5) (2.3,1.8)
    };

\end{tikzpicture}
\vspace{0.5cm}\noindent
上图，根据已知点如何预测下一个位置呢？若选用线性模型$y=a_n*x_n+a_{n-1}*x_{n-1}+\dots+a_0$可以拟合所有点，
但$y=a*x+b$简单的只是一条直线。选用哪个预测模型更好呢？更高维度的线性模型，实际上也就是多元线性回归：

\begin{equation}
y = \beta_0 + \beta_1x_1+\beta_2x_2+\dots+\beta_px_p+\epsilon
\end{equation}

\noindent
其中，$\epsilon$代表一些未知因素引入的随机误差（残差），譬如丢失一些关键特征。
这个误差会导致有些点偏离预测的函数曲线。在机器学习中，线性回归的模型描述通常使用$\hat{y}=w^Tx+b$的形式。
其中，$\hat{y}$是模型预测y应该取的值。

\begin{table}[!htbp]\centering
\begin{tabular}{|c|c|c|c|c|c|c|c|}
\hline
\multicolumn{8}{|c|}{已知数据}\\
\hline
1&2&3&4&5&6&7&8\\
\hline
(0.2,0.3)&(0.4,0.6)&(0.6,1.0)&(1.0,1.2)&(1.4,1.2)&(1.6,1.7)&(2.0,1.5)&(2.3,1.8)\\
\hline
\end{tabular}
\end{table}

对于简单的\emph{线性回归}问题，使用\emph{ND4J}和\emph{DL4J}都可以求解。
不同的是\emph{ND4J}结合最小二乘法直接公式求解，而\emph{DL4J}根据已知数据训练的结果。
\vspace{0.3cm}

1、创建神经网络
\begin{lstlisting}[language=Java]
OutputLayer outputLayer = new OutputLayer.Builder(LossFunctions.LossFunction.MSE)
                .activation(Activation.IDENTITY)
                .nIn(1)
                .nOut(1).build();

MultiLayerConfiguration conf = new NeuralNetConfiguration.Builder()
        .optimizationAlgo(OptimizationAlgorithm.STOCHASTIC_GRADIENT_DESCENT)
        .weightInit(WeightInit.ZERO)
        .updater(new Sgd(0.01))
        .list()
        .layer(0, outputLayer)
        .pretrain(false)
        .backprop(true).build();

MultiLayerNetwork net = new MultiLayerNetwork(conf);
net.init();
\end{lstlisting}

2、初始化数据
\begin{lstlisting}[language=Java]
double [] input = {
    0.2, 0.4, 0.6, 1.0, 1.4, 1.6, 2.0, 2.3
};

double [] output = {
    0.3, 0.6, 1.0, 1.2, 1.2, 1.7, 1.5, 1.8
};
INDArray inputNDArray = Nd4j.create(input, new int[]{input.length,1});
INDArray outPut = Nd4j.create(output, new int[]{output.length, 1});
\end{lstlisting}
3、机器学习过程
\begin{lstlisting}[language=Java]
for( int i=0; i<5000; i++ ) {
    net.fit(inputNDArray, outPut);
}

Map<String, INDArray> params = net.paramTable();
params.forEach((key, value) -> System.out.println("key:" + key +", value = " + value));
\end{lstlisting}

\begin{verbatim}
key:0_W, value = 0.6350
key:0_b, value = 0.4085
\end{verbatim}
\noindent
因此，拟合出的直线方程： $y=0.635*x+0.4085$

\begin{tikzpicture}[scale=1.2]
    % Draw axes
    \draw [<->,thick] (0,2.5) node (yaxis) [above] {$y$}
        |- (4,0) node (xaxis) [right] {$x$};
    
    \draw (0,0.4085) coordinate (a_1) -- (3,0.635*3+0.4085) coordinate (a_2);

    \fill[red] (0.2,0.3) circle (2pt);
    \fill[red] (0.4,0.6) circle (2pt);
    \fill[red] (0.6,1.0) circle (2pt);

    \fill[red] (1.0,1.2) circle (2pt);
    \fill[red] (1.4,1.2) circle (2pt);
    \fill[red] (1.6,1.7) circle (2pt);

    \fill[red] (2.0,1.5) circle (2pt);
    \fill[red] (2.3,1.8) circle (2pt);

    \draw[red,dashed] plot[smooth] coordinates {
        (0.2,0.3) (0.4,0.6) (0.6,1.0) (1.0,1.2)
        (1.4,1.2) (1.6,1.7) (2.0,1.5) (2.3,1.8)
    };

\end{tikzpicture}

\noindent
很显然，机器学习预测的图像，比上一页凭我直观的绘制直线更加符合实际。

\section{逻辑回归应用}
\section{多分类问题}
\section{聚类问题}
\section{卷积神经网络}
\section{循环卷积神经网络}







 %AI
% !Mode:: "TeX:UTF-8"
\part{综合运用}
\label{part:ai_exercise}

\newpage

%二维码
\begin{center}{\includegraphics[width=3cm]{barcode.png}}\end{center}\par

十年前，\emph{手写识别}还只是专业技术人员的玩物，而如今\emph{手写识别}已经变成深度学习的入门教程。
而\emph{模式识别}（Pattern Recognition）意味着需要专业人员提供很多特征。
模式识别从19世纪50年代开始出现，在20世纪80年代左右曾经风靡一时，也是信息科学和人工智能的重要组成部分。
它被应用于图像分析与处理、语音识别、数据挖掘等很多方面。
尽管模式识别也很高大上，也有了较长时间的应用，但效果一直差强人意。

这种由人提取特征交给机器，再交给机器去判断的方法识别率一直差强人意。
而\emph{机器学习}可以从海量数据学习\emph{特征}，使用已知经验数据（样本）自己提炼特征。
另外，再强调一下概念，深度学习只是机器学习算法的一种，而人工智能也不限于机器学习一种。

\begin{figure}[!htb] \centering
\begin{tikzpicture}
\node(learning) [draw, rectangle] at (2,0) {学习};
\node(knowledge) [draw, ellipse] at(2,1.5) {经验};
\draw[->](0,0) node[left]{新数据} -- node[above]{输入} (learning);
\draw[->](learning) -- node[above]{预测} (4,0) node[right]{人工智能};
\draw[->](knowledge) -- (learning);
\end{tikzpicture}
\end{figure}

大数据和深度学习，让我们能够摆脱人工特征提取的繁琐和束缚，让更多的人参与人工智能并为社会创造价值。
互联网+没有过时，再加上分布式和云计算，使得\emph{人工智能+}如鱼得水，再次迎来十年兴盛，相信接下来是数据为王的时代！

\chapter{大数据和分布式}
\label{chap:java_ex_cloud}

\section{云计算}

云计算是基于互联网的一种全新的模式。通过“云”我们可以轻松地存储数据，运行应用，
甚至进行强度很大的运算。云计算实质上实现了对计算机资源的完全掌控，
我们可以很轻易地通过云对资源进行分配，例如分配内存、处理器等计算资源。

机器学习和云计算、大数据的关系十分密切。大数据技术具有大量的数据与资源，能促进
机器学习的进一步发展。反过来，机器学习的发展也能促进数据挖掘，分析等方面
的进步。

\section{边缘计算}
边缘计算的定义是在靠近数据生成的设备端，就近提供服务。相比于云计算需要
将所有数据上传至云端的特性，边缘计算具有天然的优势。第一，由于边缘计算强调
“边缘”，数据的处理相比于云计算更为实时，高效；第二，边缘计算处理的数据多为
“小数据”，在数据存储和计算上成本都比较低；第三，边缘计算的出现降低了
对网络带宽的需求。随着物联网的不断发展，越来越多的联网设备将出现在
我们的生活中，这大大增加了网络传输的压力。边缘计算则通过数据本地处理
解决了这一问题。

边缘计算是对云计算的一种补充与完善。思科预计到2022年，移动设备和连接数量将达到123亿
\footnote{数据来源思科官网《思科移动网络VNI预测（2017-2022 年）》}。如果仅仅用
云计算来解决这么多互联设备的数据量往往会造成资源浪费，及时性较差，以及隐私泄露
等问题。边缘计算在一方面能减轻云端的负担，另一方面也能保护用户的隐私，具有重要
的实际意义。

\section{开源架构}
随着云计算和大数据的不断发展，人们对于数据存储和数据处理的效率
提出了越来越严格的要求，因此一批优秀的开源架构应运而生。

\subsection{Hadoop}
Hadoop是基于分布式计算的大数据开源架构，它的特点在于能高效，可靠地对
大量数据进行处理。Hadoop由HDFS和MapReduce组成，分别负责数据的分布式
储存和分布式计算。

HDFS提供了大规模存储的分布式解决方案，可以横跨成千上百台机器，操作却像对本地文件
进行操作一样简单，并能存储TB甚至PB级别的特大文件。这些优势来源于在HDFS中，
文件的储存是以较大的抽象块为单位，这样做的好处在于能将文件存储于
不同的硬盘中，使得文件的大小不受限于硬盘大小；同时较大的抽象块有利于对大型
文件进行寻址操作，最小化寻址开销。

除了抽象块，HDFS还包含了NameNode和DataNode。用户在用HDFS进行存储时，客户端首先
会将文件切块，再由DataNode储存在多个硬盘中。同时，NameNode负责记录每一个
文件的切块信息，以及每一块具体的存储机器。这样的系统就构成了NameNode为主服务器，多个
DataNode为从服务器的HDFS文件系统。

MapReduce是Hadoop的计算框架，能够以容错、可靠的方式并行处理海量的数据。
MapReduce主要包括了两个主要部分：Map阶段和Reduce阶段。



\subsection{Storm}
Storm是类似于Hadoop的实时数据处理框架，两者的区别在于Storm用于数据
的实时计算，Hadoop用于对离线数据进行处理。同时，Storm将数据存储在
内存中，这也有区别于Hadoop的HDFS存储方式。Storm也提供了类似于MapReduce
的简单编程模型，便于进行开发。

一个完整的Storm集群应该包括Nimbus，Zookeeper和Supervisor，工作架构
如图\ref{storm_framework}所示。其中，Nimbus负责接受客户端传来的拓扑代码
并分拆task，Zookeeper负责Nimbus和Supervisor之间的通信和协调，Supervisor
接收并处理数据。这样高效的工作架构使得Storm具有实时处理大批量数据的能力，
同时具有高扩展性，可以通过增加节点实现性能的线性提高。

~~~~
\begin{figure}[!ht]
\begin{center}
\begin{tikzpicture}
\tikzstyle{arrow} = [thick, -, >= stealth]
\tikzstyle{node} = [rectangle, rounded corners, minimum width = 2cm, minimum height=1cm,text centered, draw = black]
\node[node](Nimbus){Nimbus};
\node[node, right of = Nimbus, xshift = 2cm](zookeeper2){zookeeper};
\node[node, below of = zookeeper2, yshift = -0.5cm](zookeeper3){zookeeper};
\node[node, below of = zookeeper2, yshift = 2.5cm](zookeeper1){zookeeper};
\node[node, right of = zookeeper1, xshift = 2cm](supervisor1){supervisor};
\node[node, right of = zookeeper2, xshift = 2cm](supervisor2){supervisor};
\node[node, right of = zookeeper3, xshift = 2cm](supervisor3){supervisor};
\draw[arrow] (Nimbus)--(zookeeper1);
\draw[arrow] (Nimbus)--(zookeeper2);
\draw[arrow] (Nimbus)--(zookeeper3);
\draw[arrow] (zookeeper1)--(supervisor1);
\draw[arrow] (zookeeper1)--(supervisor2);
\draw[arrow] (zookeeper1)--(supervisor3);
\draw[arrow] (zookeeper2)--(supervisor1);
\draw[arrow] (zookeeper2)--(supervisor2);
\draw[arrow] (zookeeper2)--(supervisor3);
\draw[arrow] (zookeeper3)--(supervisor1);
\draw[arrow] (zookeeper3)--(supervisor2);
\draw[arrow] (zookeeper3)--(supervisor3);
\end{tikzpicture}
\caption{Storm架构示意图}
\label{storm_framework}
\end{center}
\end{figure}

\subsection{Spark}
Spark是一种通用的大数据计算框架，与Hadoop的MapReduce，Storm的
流式实时计算引擎十分相似。同时Spark也提供了大数据和机器学习方面的
框架，逐渐形成了大数据处理的一站式解决平台。

Spark的特点是支持大型，低延迟的数据处理方式。拓展了常用的MapReduce计算
模型。由于Spark支持在内存中进行计算，在必要时依赖硬盘进行复杂运算，因此
Spark比MapReduce更高效。

Spark是MapReduce的一种替代方案，兼容HDFS，Hive等存储方案，能够方便
地融入Hadoop的生态环境，以弥补MapReduce的不足。

\section{示例：人脸识别}



\chapter{车牌识别}
\label{chap:java_ex_recongnition}


手写识别机器学习领域中的一个经典问题。自从深度学习的兴起，MNIST已经是一个被”嚼烂”了的数据集。
该问题要解决的是把28x28像素的灰度手写数字图片，识别为0到9的数字。
而\emph{手写算式识别}在数字识别的基础上，需要识别运算符号，并输出计算结果。

\begin{figure}[!htb]
\centerline{\includegraphics[width=.1\figwidth]{images/handwriting_math.png} $\to$ 识别: $3+5 \times 6$} 
\label{fig:part4_handwriting_math}
\caption{手写算式识别}
\end{figure}

限于篇幅和笔者能力，本演示代码仅支持四则运算（$+ - \times \div$），其它运算符号留给同学们研究。
利用深度学习算法，为用户提供手写算式识别，可以极大的提高用户体验。
目前，有很多成熟的产品可以借鉴，相关的论文也非常多。

手写数字识别属于图像分类，可用的算法有：KNN、SVM、BP和CNN。
其中KNN和SVM比较善于数据挖掘，用做手写数字识别的话，比较难达到实用目的；
而BP使用全连接，难于构建深度网络，并存在梯度消失和过拟合问题。
本章例子，选用CNN深度学习训练手写算式识别。

\begin{figure}[!htb] \centering 
\begin{tikzpicture}
\node at (0,0) {\includegraphics[width=5cm,height=2cm]{images/handwriting_math.png}};
\draw (-2.5,1) rectangle (-1.6,-1);
\draw (-1.6,1) rectangle (-0.8,-1);
\draw (-0.8,1) rectangle (0.3,-1);
\draw (0.3,1) rectangle (1.3,-1);
\draw (1.3,1) rectangle (2.4,-1);
\end{tikzpicture}
\end{figure}

现在我们要识别一个算式，其实和MINIST数字识别没有太多差别，主要是增加了图像切割和四则运算符的识别任务。
为了降低难度，约束字符之间最好有一些间距，不要字符连笔。


\section{图像切割}
手写算式的输入的形式，有2种：拍照、手绘。照片要经过图像处理之后，才能进行切割；
而手绘可直接采集像素轨迹，就省掉了这一步。对于照片的预处理，需要先二值化和降噪把手写数字区域保留下来。
如果考虑图片的背景干扰，可先使用Canny算子检测边缘，这借助OpenCV的API很容易实现。

我们这种要求手写数字不黏连的情况，把Canny输出的图像往X轴投影，就能得出每个数字的横向宽度坐标。
虽然要求数字不黏连，但难保有横向重叠的情况，必须用类似\emph{范水法}进行切割。

\begin{figure}[!htb] \centering 
\begin{tikzpicture}
\node at (0,0) {\includegraphics[width=5cm]{images/handwriting_input.png}};
\node at (6,0) {\includegraphics[width=5cm]{images/handwriting_out.png}};
\end{tikzpicture}
\caption{图像处理：二值化和边缘识别(OpenCV)}
\end{figure}

\noindent
这种图像预处理，OpenCV\footnote{OpenCV是一个开源的跨平台视觉库，实现了图像处理和计算机视觉方面的很多通用算法，
支持Linux、Windows、Android和Mac OS操作系统上。
它轻量且高效，同时提供C++、Python、MATLAB等语言的接口。}非常擅长，不需要使用机器学习算法。


\noindent
\emph{Java代码：}
\begin{lstlisting}[language=Java]
System.loadLibrary(Core.NATIVE_LIBRARY_NAME);
Mat img = Imgcodecs.imread(src);
Imgproc.cvtColor(img, img, Imgproc.COLOR_BGR2GRAY);
Imgproc.Canny(img, img, threshold, threshold * 3, 3, true);
Imgcodecs.imwrite(dstImg, img);
\end{lstlisting}

\noindent
\emph{Python代码：}
\begin{lstlisting}[language=Python]
import cv2
from matplotlib import pyplot as plt

img = cv2.imread("handwriting_input.png", 0)
edges = cv2.Canny(img, 30, 60)
plt.subplot(1,2,1), plt.imshow(img, cmap='gray')
plt.title('Original Image'), plt.xticks([]), plt.yticks([])
plt.subplot(1,2,2), plt.imshow(edges, cmap='gray')
plt.title('Edge Image'), plt.xticks([]), plt.yticks([])
plt.show()
\end{lstlisting}

\noindent
实现Canny算子不在本书讨论范围，GitHub上的一份Java实现接近OpenCV的效果
\href{https://github.com/enzoftware/images-processing/blob/master/algorithms/Canny%20Filter/src/com/company/CannyFilter.java}
{\srclink{Github}上获取。}而CSDN上的一份实现简化了很多处理细节，有一定差距。有条件的同学，建议使用OpenCV的接口。
遇到的OpenCV问题，也容易在网络上检索。
\vspace{0.3cm}
\hrule
\begin{multicols}{2}
\begin{itemize}
\item[1.]转换为灰度图像
\item[2.]高斯模糊处理
\item[3.]根据梯度计算图像边缘幅值与角度
\item[4.]非最大信号压制处理（边缘细化）
\item[5.]双阈值边缘连接处理
\item[6.]二值化图像输出结果
\end{itemize}
\includegraphics[width=.12\figwidth]{images/Lenna.png}
\includegraphics[width=.12\figwidth]{images/Lenna_out.png}
% \end{tabular}
\end{multicols}
\hrule
\ \\
经过Canny算子处理之后，切割图像就非常容易了。把Canny输出的二值图像，在水平方向上投影，对数字进行竖向分割；
同理，在纵向再做一次，就能切割得紧凑一些。以下，是在原图上切割的效果 ：

\begin{figure}[!htb] \centering 
\begin{tikzpicture}
\node at (0,0) {\includegraphics[width=1cm]{images/3.png}};
\node at (2,0) {\includegraphics[width=1cm]{images/plus.png}};
\node at (4,0) {\includegraphics[width=1cm]{images/5.png}};
\node at (6,0) {\includegraphics[width=1cm]{images/times.png}};
\node at (8,0) {\includegraphics[width=1cm]{images/6.png}};
\end{tikzpicture}
\end{figure}

\noindent
Minist数据集是白底黑字的，使用它训练出来的模型，就很难识别我们的手写算式。所以，有必要在二值图像上切割。
\begin{figure}[!htb] \centering 
\begin{tikzpicture}
\node at (0,0) {\includegraphics[width=1cm]{images/3_c.png}};
\node at (2,0) {\includegraphics[width=1cm]{images/plus_c.png}};
\node at (4,0) {\includegraphics[width=1cm]{images/5_c.png}};
\node at (6,0) {\includegraphics[width=1cm]{images/times_c.png}};
\node at (8,0) {\includegraphics[width=1cm]{images/6_c.png}};
\end{tikzpicture}
\end{figure}

\noindent
但Canny生成的是空心字，并且有很多断线，使用FloodFill算法是填不满的。
笔者采用的是膨胀法，对于每一个\emph{白色点}向周围扩张6个像素，这样的结果还不错，差不多和Minist数据集很接近了。
\begin{figure}[!htb] \centering 
\begin{tikzpicture}
\node at (0,0) {\includegraphics[width=1cm]{images/3_b.png}};
\node at (2,0) {\includegraphics[width=1cm]{images/plus_b.png}};
\node at (4,0) {\includegraphics[width=1cm]{images/5_b.png}};
\node at (6,0) {\includegraphics[width=1cm]{images/times_b.png}};
\node at (8,0) {\includegraphics[width=1cm]{images/6_b.png}};
\end{tikzpicture}
\end{figure}

\noindent
尺寸也要标准化，做成和Minister一样的$28\times28$，这个对图片缩放就可以实现。


\section{算式识别}


\section{结果评估}
\chapter{写诗机器人}
\label{chap:java_ex_poem}

\section{word2vec介绍}
word2vec是google在2013年推出的一个NLP工具,它的特点是将所有的词向量化,这样词与词之间就可以定量的去度量他们之间的关系,挖掘词之间的联系。

\section{如何写诗}


 %Ex.

\backmatter
\bookmarksetup{startatroot}

\appendix

\small{
\bibliography{dl4j}
\bibliographystyle{natbib}
\clearpage
}

\printglossary[title=术语]

\end{document}